\documentclass[12pt]{article}
\usepackage[utf8]{inputenc}
\usepackage[T1]{fontenc}
\usepackage{hyperref}
\title{Mixtec, Tamazola Language Resources for K-12 Curriculum Development}
\author{Compiled by Claude Sonnet 5 (claude-sonnet-5) for Andr\'as Kornai}
\date{}
\begin{document}
\maketitle

\section{Introduction}
Tamazola Mixtec (ISO 639-3 vmx; Glottolog tama1339) is an Otomanguean language of Oaxaca, Mexico. The searches found little. No textbooks, primers or school programmes appeared.

\section{Primary Educational Resources}
None were found.

\section{Linguistic and Reference Materials}
\begin{itemize}
\item AILLA (UT Austin) syntax elicitation recordings from Santiago Tamazola, 1979, with PDFs. Researcher: Kathryn Josserand: \url{https://islandora-ailla.lib.utexas.edu/islandora/object/ailla%3A129150}
\item A second AILLA item from the same collection: \url{https://islandora-ailla.lib.utexas.edu/islandora/object/ailla%3A129158}
\item Glottolog entry: \url{https://glottolog.org/resource/languoid/id/tama1339}
\end{itemize}

\section{Religious and Community Resources}
\begin{itemize}
\item Global Recordings Network audio Bible lessons, 40 sections: \url{https://globalrecordings.net/en/language/vmx}
\item Joshua Project people-group profile, which lists the written and Bible translation status as unknown: \url{https://joshuaproject.net/people_groups/14700}
\end{itemize}

\section{Digital and Audio Resources}
The AILLA recordings and the Global Recordings audio above are the only audio found.

\section{Curriculum Development Context}
Existing material is archival or evangelistic. A curriculum would need new development with Tamazola community speakers and regional Mexican indigenous education bodies. None of those were located in the searches.
\end{document}